\documentclass[
10pt, % Main document font size
a4paper, % Paper type, use 'letterpaper' for US Letter paper
oneside, % One page layout (no page indentation)
%twoside, % Two page layout (page indentation for binding and different headers)
headinclude,footinclude, % Extra spacing for the header and footer
BCOR5mm, % Binding correction
]{scrartcl}
%%%%%%%%%%%%%%%%%%%%%%%%%%%%%%%%%%%%%%%%%
% Arsclassica Article
% Structure Specification File
%
% This file has been downloaded from:
% http://www.LaTeXTemplates.com
%
% Original author:
% Lorenzo Pantieri (http://www.lorenzopantieri.net) with extensive modifications by:
% Vel (vel@latextemplates.com)
%
% License:
% CC BY-NC-SA 3.0 (http://creativecommons.org/licenses/by-nc-sa/3.0/)
%
%%%%%%%%%%%%%%%%%%%%%%%%%%%%%%%%%%%%%%%%%

%----------------------------------------------------------------------------------------
%	REQUIRED PACKAGES
%----------------------------------------------------------------------------------------

\usepackage[
nochapters, % Turn off chapters since this is an article        
beramono, % Use the Bera Mono font for monospaced text (\texttt)
eulermath,% Use the Euler font for mathematics
pdfspacing, % Makes use of pdftex’ letter spacing capabilities via the microtype package
dottedtoc % Dotted lines leading to the page numbers in the table of contents
]{classicthesis} % The layout is based on the Classic Thesis style

\usepackage{arsclassica} % Modifies the Classic Thesis package
\usepackage{ulem}
\usepackage{graphicx}
\usepackage{float}
\usepackage{algorithm2e}
\usepackage[export]{adjustbox}
\usepackage[utf8]{inputenc}  
\usepackage[english,frenchb]{babel} 
\usepackage[T1]{fontenc} % Required for accented characters
\usepackage{multicol}
\usepackage{graphicx} % Required for including images
\graphicspath{{Figures/}} % Set the default folder for images

\usepackage{enumitem} % Required for manipulating the whitespace between and within lists

\usepackage{lipsum} % Used for inserting dummy 'Lorem ipsum' text into the template


\expandafter\def\csname ver@subfig.sty\endcsname{}


\usepackage{svg}
\usepackage{hyperref}

\usepackage{amsmath,amssymb,amsthm} % For including math equations, theorems, symbols, etc

\usepackage{varioref} % More descriptive referencing

%----------------------------------------------------------------------------------------
%	THEOREM STYLES
%---------------------------------------------------------------------------------------

\theoremstyle{definition} % Define theorem styles here based on the definition style (used for definitions and examples)
\newtheorem{definition}{Definition}

\newtheorem*{remark}{Remark}
\theoremstyle{plain} % Define theorem styles here based on the plain style (used for theorems, lemmas, propositions)
\newtheorem{theorem}{Theorème}
\newtheorem{lemma}[theorem]{Lemme}
\newtheorem{proposition}[theorem]{Proposition}

\theoremstyle{remark} % Define theorem styles here based on the remark style (used for remarks and notes)

%----------------------------------------------------------------------------------------
%	HYPERLINKS
%---------------------------------------------------------------------------------------

\hypersetup{
%draft, % Uncomment to remove all links (useful for printing in black and white)
colorlinks=true, breaklinks=true, bookmarks=true,bookmarksnumbered,
urlcolor=webbrown, linkcolor=RoyalBlue, citecolor=webgreen, % Link colors
pdftitle={}, % PDF title
pdfauthor={\textcopyright}, % PDF Author
pdfsubject={}, % PDF Subject
pdfkeywords={}, % PDF Keywords
pdfcreator={pdfLaTeX}, % PDF Creator
pdfproducer={LaTeX with hyperref and ClassicThesis} % PDF producer
}

\hyphenation{Fortran hy-phen-ation} 

\title{\normalfont{FTML Exercices 2 solutions}\\
} 
\author{\spacedlowsmallcaps{}
} 
\date{\today}
\begin{document}
\renewcommand{\sectionmark}[1]{\markright{\spacedlowsmallcaps{#1}}}
\lehead{\mbox{\llap{\small\thepage\kern1em\color{halfgray} \vline}\color{halfgray}\hspace{0.5em}\rightmark\hfil}} 
\pagestyle{scrheadings}
\maketitle 
\setcounter{tocdepth}{3} 

\tableofcontents

\section{\large\color{Blue}Ordinary least squares}

\subsubsection{Solution}%
\label{ssub:Solution}

\textbf{1)} On connaît déjà le gradient de l'application $f: x\mapsto ||x||^2$, qui vaut
$2x$.

Si on considère l'application $r$:

$$
r = \left\{
    \begin{array}{ll}
        \mathbb{R}^d \rightarrow \mathbb{R}^n  \\
	\theta\mapsto X\theta-y
    \end{array}
\right.
$$

alors $g = f \circ r$. Comme tout est différentiable, on en déduit qu'en notant
$L$ les jacobiennes:

\begin{equation}
    L_{\theta}g = L_{X\theta-y}fL_{\theta}r
\end{equation}

ou bien on considérant le gradient (qui est la transposée de la jacobienne
quand l'application est à valeurs dans $ \mathbb{R} $):

\begin{equation}
    \nabla_{\theta}g(\theta) = (L_{\theta}r)^T\nabla_{x}f(X\theta-y)
\end{equation}

Or $L_{\theta}r = X$. Donc

\begin{equation}
    \nabla_{\theta}g(\theta) = 2X^T(X\theta-y)
\end{equation}



\textbf{2a)} Soit $\alpha\in [0, 1]$, et $x, y\in \mathbb{R}^d$.

\begin{equation}
    \begin{aligned}
        \label{eq:}
	(f\circ s)(\alpha x+(1-\alpha)y)&= f(s(\alpha x+(1-\alpha)y))\\
	&= f(\alpha s(x)+(1-\alpha)s(y))\\
	&\leq \alpha f(s(x))+(1-\alpha)f(s(y))
    \end{aligned}
\end{equation}

\textbf{2b)} Soit $\alpha\in [0, 1]$, et $x, y\in \mathbb{R}^n$.

\begin{equation}
    \begin{aligned}
        \label{eq:}
	||\alpha x+(1-\alpha)y||&\leq ||\alpha x||+||(1-\alpha)y||\\
	&=\alpha ||x||+(1-\alpha )||y||
    \end{aligned}
\end{equation}

\textbf{2c)} Soit $\alpha\in [0, 1]$, et $x, y\in \mathbb{R}$.

Since $a$ is convex,

\begin{equation}
	 a(\alpha x+(1-\alpha)y)\leq \alpha a(x)+(1-\alpha)a(y)
\end{equation}

Since $w$ is increasing,

\begin{equation}
    w(a(\alpha x+(1-\alpha)y))\leq w(\alpha a(x)+(1-\alpha)a(y))
\end{equation}

Since $w$ is convex,

\begin{equation}
    w(\alpha a(x)+(1-\alpha)a(y))\leq \alpha(w(a(x))+(1-\alpha)w(a(y))
\end{equation}

Finally,

\begin{equation}
    (w\circ a)(\alpha x+(1-\alpha)y))\leq \alpha(w\circ a)(x))+(1-\alpha)(w\circ a)(y)
\end{equation}

\textbf{2d)} Soit $\alpha\in [0, 1]$, et $x, y\in \mathbb{R}^n$.
\begin{equation}
	\begin{aligned}
		\label{eq:}
		f(\alpha x + (1-\alpha) y) &= u(\alpha x + (1-\alpha) y + \beta)\\
		&= u(\alpha x + (1-\alpha) y + (\alpha + 1-\alpha)\beta)\\
		&= u(\alpha x + (1-\alpha) y + \alpha\beta + (1-\alpha)\beta)\\
		&= u(\alpha (x+\beta) + (1-\alpha) (y+\beta))\\
		&\leq \alpha u(x+\beta) + (1-\alpha)u(y+\beta))\\
		&= \alpha f(x) + (1-\alpha)f(y)\\
	\end{aligned}
\end{equation}

\textbf{2e)} On utilise:
\begin{itemize}
    \item le point \textbf{c)} avec $w: t\mapsto t^2$ et $a$ l'application
        norme sur $\mathbb{R}^n$ pour montrer que $u:x\mapsto ||x||^2$ est convexe.
    \item le point \textbf{d)} avec $\beta=-y$ pour montrer que $f:x\mapsto ||x-y||^2$ est convexe de $\mathbb{R}^n$ dans $\mathbb{R}$.
    \item le point \textbf{a)} appliqué à $g=f\circ s$ avec $s:\theta \mapsto X\theta$ linéaire de $ \mathbb{R}^d$ dans $\mathbb{R}$.
\end{itemize}

\textbf{3)} Puisque la fonction objectif est convexe, et différentiable, résoudre l'annulation du gradient est équivalent à trouver les minima locaux. Si la matrice $X^TX$ est inversible, alors l'annulation du gradient a bien une unique solution donnée par la formule OLS.

\end{document}
